The results chapter reported each of the five dimensions as an independent release-wise series. This chapter draws those series together in answer to the research questions, relates the synthesis to the prior work reviewed in \autoref{ch:relatedwork}, states the limitations that bound the findings, and identifies directions for future work. Consistent with the exploratory-descriptive framing established in \autoref{ch:introduction}, the synthesis below characterises and situates what was measured; it does not assert a causal or correlational relationship between the dimensions, and the inferential analysis that such claims would require remains a direction for future work.

\section{Principal Findings}
\label{sec:principal-findings}

The central observation that the primary research question yields is a divergence between the size of the persistent core and the complexity of its constituent parts. Over the 25 years from Potato to Trixie, the persistent core grew by roughly 5.3-fold in lines of code, from approximately 0.92 to approximately 4.92 million (\autoref{fig:nloc}), and its installed footprint grew roughly fourfold, from approximately 20 to approximately 82~MB (\autoref{fig:installed-size}), for a population that is fixed by construction. Across the same span, however, the per-function cyclomatic complexity did not rise: the median held at exactly 3 in every one of the 13 releases, and the mean rose by only 0.29 to a peak of 6.58 at Jessie before declining to 6.08 at Trixie, a value below its starting point of 6.29 (\autoref{fig:complexity}). The per-package dependency count was similarly stationary, with a median of one edge per package in every release after Potato (\autoref{fig:dependency-distribution}). The descriptive picture is therefore one of growth by accretion: as the number of analysed functions more than tripled, from approximately 15{,}900 to approximately 55{,}000, the core expanded through the addition of many small, low-complexity, loosely coupled units rather than through an increase in the complexity or coupling of the typical unit. The fifth dimension that the primary question names, distribution-specific patch volume, is synthesised below together with the secondary research question that sharpens it.

The secondary research question on composition is answered by a non-monotonic package count whose movements are attributable to packaging structure rather than to changes in capability. The minbase count rises and falls between 44 and 95 packages across the releases (\autoref{fig:package-count}), and the two largest movements both resolve into composition effects on inspection: the peak at Jessie is produced by the simultaneous presence of legacy, newly arrived, and transient packages during the \texttt{systemd} init-system transition (\autoref{fig:init-transition}), and the subsequent decline is produced by library soname bumps, the 64-bit \texttt{time\_t} rebuild, the departure of dependency chains, and the requalification of what a minimal chroot requires. Because splitting and renaming packages move the count without adding or removing code, the count is a weak indicator of size; the installed-size and lines-of-code dimensions are the stronger signals. Those stronger signals indicate that recent growth is carried by the transient set rather than by the long-lived core: the persistent footprint plateaus near 82~MB from Buster onward while the gap to the full minbase widens to approximately 30 to 40~MB.

The secondary research question on distribution-specific patches is answered only within the window in which the metric is comparable. As established in \autoref{sec:patch-inconsistency}, patch-line counts under the \texttt{1.0} source format are inflated by non-code material bundled into a single diff, so the series is interpreted only from Jessie onward under the \texttt{3.0 (quilt)} format. Within that window the total patch volume across the persistent core declines from approximately 663{,}000 to approximately 51{,}000 lines (\autoref{fig:patch-size}), a total dominated in every release by \texttt{glibc}, which alone accounts for approximately 570{,}000 of the Jessie figure and approximately 37{,}000 of the Trixie figure. The file-level structure changed in parallel, from a single bundled \texttt{.diff.gz} to many named patches under the DEP-3 convention~\cite{dep3}; \texttt{glibc}, for instance, moves from one diff to 281 named patches at Jessie. The decline therefore describes a contraction in the bulk of Debian's named modifications to the persistent core, though, because the metric counts patch lines including diff context rather than semantic changes, it indicates the bulk of the distribution's modifications rather than the effort required to produce them.

One structural event recurs across these dimensions. The arrival of \texttt{systemd} in the required set at Jessie is visible simultaneously as the peak in package count, the largest single step in the minbase installed-size series, and the single dependency outlier of 17 edges that is the highest value in the dataset. This co-visibility is a descriptive coincidence: a single packaging decision left a trace in three metrics computed independently from the same release. It does not show that any one dimension drives another, a claim the descriptive design of this thesis cannot establish.

\section{Relation to Prior Work}
\label{sec:relation-prior-work}

The continuing growth of the persistent core is a direct instance of the tendency that Lehman formalised and that Wirth restated. Lehman's laws hold that actively used software continually grows unless effort is invested otherwise~\cite{Lehman1980Laws}, and Wirth argued that this growth had become a structural problem in which resource demands outpaced functional gains~\cite{Wirth1995Plea}. The 5.3-fold growth in lines of code of a fixed package population, observed here for a minimal base system rather than for an application, is consistent with both accounts. The companion proposition that complexity increases unless deliberately countered holds in the present data only at the aggregate level: the total complexity summed over all functions rises with the function count, and the maximum per-function complexity grows from 273 to a peak of 1{,}125 (\autoref{sec:results-cc}), yet the complexity of the typical function does not increase. The data therefore supports the continuing-growth tendency while qualifying the increasing-complexity tendency at the granularity of the individual function.

A particularly close correspondence is with the kernel study of Israeli and Feitelson. They reported that the total complexity of the Linux kernel rises over time while the average complexity of individual functions falls, because new code is introduced as many small, low-complexity functions~\cite{israeli2010linux}. The persistent Debian core reproduces this mechanism on a different corpus: the function count triples, the median per-function complexity is stationary, and the mean declines after its Jessie peak. The agreement extends their kernel-specific observation to the base of a complete distribution and supports their methodological argument that total and average complexity must be reported together, because either measure in isolation would misrepresent the trend.

The growth-rate studies of the kernel admit a more guarded comparison. Godfrey and Tu reported super-linear growth in kernel lines of code~\cite{godfrey2000linux}, and Koren contrasted the faster accumulation of code in development branches against stable ones~\cite{koren2006linux}. The persistent core analysed here grows smoothly and monotonically, but over a package population that is fixed by construction, so its trajectory measures within-package growth and excludes the arrival of new modules that drives super-linear whole-system curves. The present design isolates a different component of growth rather than contradicting these results. At the broader level of empirical validation, Neamtiu et al.\ and Vasa each found continuing growth and continuing change to be the most robustly confirmed of Lehman's laws across many individual projects~\cite{neamtiu2009towards,vasa2010growth}; the present work adds an observation of the same tendency at the granularity of a whole distribution's base.

The dependency findings relate to the network studies of the Debian ecosystem. LaBelle and Wallingford, De Sousa et al., and Horváth each characterised Debian's inter-package dependency graph as scale-free, with new edges accumulating by preferential attachment~\cite{interPackageDependency,sousa2009debian,horvath2012linuxdeps}. The present result, in which the median package has declared between zero and two dependencies for the full 25-year span while all movement is confined to a small number of hub packages such as \texttt{systemd}, is consistent with a hub-dominated structure of that kind. The agreement is suggestive rather than a replication, because the measurement here covers only the small minbase subgraph rather than the full archive of more than eighteen thousand packages that those studies analysed.

The patch-volume dimension connects to the literature on distribution-level maintenance. Becker et al.\ analysed Debian maintainer scripts at scale with the CoLiS platform~\cite{becker2022colis}, and Lin et al.\ quantified the share of fixes contributed locally by distribution maintainers rather than propagated from upstream~\cite{lin2022upstream,lin2023vulnerability}. The contraction of the named patch set in the quilt era describes one facet of this maintenance work for the persistent core. The metric here counts patch lines including diff context rather than semantic changes, however, and CoLiS examined maintainer scripts rather than source patches, so the present result complements these studies on a distinct artefact without measuring maintenance effort directly. Molnar and Motogna observed that distribution-level metric trends can diverge from the trends of individual projects~\cite{molnar2020longitudinal}; the present aggregates echo this, in that they are dominated by single packages, with \texttt{glibc} governing both the installed-size and patch-volume totals and a single \texttt{perl} source file governing the maximum complexity. Ronchieri et al.\ established a cyclomatic-complexity baseline for a large scientific toolkit to characterise its maintainability~\cite{ronchieri2016software}; the per-function values reported here, a median of 3 and a mean near 6, provide a comparable baseline for the base of a general-purpose distribution.

Finally, the findings bear on the motivation set out in the introduction. The association that Nguyen et al.\ reported between rising complexity and increased maintenance cost and defect density~\cite{nguyen2017impact}, together with the emphasis that Sherlock et al.\ place on users being able to read and understand the software they run~\cite{Sherlock2018OpenSS}, framed complexity growth as a concern for a base system of Debian's reach. The stationary per-function complexity observed here is consistent with a core whose individual units did not become locally harder to read or maintain over 25 years, even as the corpus grew fivefold. This observation is consistent with those concerns; it is not a causal claim about maintenance outcomes, which the present data cannot support.

\section{Limitations}
\label{sec:limitations}

Cyclomatic complexity is measured only for C and C++ source, because \texttt{pmccabe} supports no other languages. The complexity results therefore describe the compiled C and C++ core of the base system and do not capture packages written in interpreted languages. The per-release complexity aggregates are consequently computed over 23 of the 25 persistent packages: \texttt{debconf} and \texttt{base-files} contribute to the lines-of-code totals but contain no C or C++ source and are therefore absent from the complexity figures. This restriction does not bias the trend within the C and C++ code, but the complexity findings should not be read as a statement about the complexity of the persistent core in its entirety.

The patch-volume metric carries two further measurement limitations. The first, established in \autoref{sec:patch-inconsistency}, is that counts are comparable only within the \texttt{3.0 (quilt)} era from Jessie onward, because the \texttt{1.0} format bundles changelogs, documentation, and other non-code material into a single diff and so inflates the earlier counts; the metric consequently describes Debian's code-level modifications for fewer than half of the releases studied. The second is that the metric counts physical patch lines, including unchanged context lines, rather than the number of semantic changes a patch makes, so it indicates the bulk of the distribution's modifications and not the effort required to produce or maintain them. The lines-of-code metric shares the character of the first concern: \texttt{tokei} counts physical lines in the upstream tarballs, which include generated code, vendored third-party code, and test suites, and the composition of a tarball can vary between releases independently of the code that the package actually builds.

A construct limitation applies to the use of cyclomatic complexity as a measure of complexity in general. Cyclomatic complexity captures the control-flow structure of individual functions and does not measure inter-procedural coupling, data-flow complexity, architectural complexity, or the cognitive difficulty of reading code. The stationary per-function complexity reported here should therefore be understood as a statement about control-flow structure at the function level, and not as evidence that the core has not become more complex along dimensions the metric does not observe.

Two limitations bound the scope to which the findings generalise. First, the persistent core is defined by survival: only the source packages present in the minbase of at least 11 of the 13 releases from Potato to Trixie qualify. This construction yields a conservative and deliberately stable subset, and the long-lived foundational packages it selects need not be representative of the base system as a whole, still less of a full Debian installation. Second, the minbase variant is a minimal, chroot-oriented installation rather than a desktop or server system, so the trajectories reported here characterise the foundation on which larger installations are built and do not describe the size or complexity that a user-facing system would exhibit. The analysis also covers a single distribution, and the processor architecture on which it is measured changes once within the window: releases up to Lenny are measured on i386 and releases from Squeeze onward on amd64 (\autoref{tab:releases}). The installed-size metric depends on the architecture for which a package is built, so the size series crosses this boundary at Squeeze, whereas the lines-of-code, complexity, and patch metrics are computed from architecture-independent source archives and are unaffected. The restriction to one distribution further means that the observed trends cannot be separated into those specific to Debian and those common to long-lived Linux distributions in general.

A final limitation concerns the provenance of the data. The early releases are reconstructed from the Debian snapshot archive~\cite{debianSnapshot}, whose historical source packages occasionally required repacking, and the diagnosis in \autoref{sec:webapp} showed that defects in the extraction of these archives produced measurement artefacts in early pipeline runs. Although the artefacts identified there were corrected, the oldest releases rest on the least uniform source material, and undetected irregularities in that material would affect the earliest points of each series most.

\section{Future Work}
\label{sec:future-work}

The limitations above indicate the most immediate extensions. Replacing or supplementing \texttt{pmccabe} with complexity tools for interpreted languages would bring \texttt{debconf}, \texttt{base-files}, and the growing \texttt{perl} share of the corpus within the scope of the complexity analysis, closing the gap between the lines-of-code and complexity populations. A semantic analysis of the patch archive, counting the changes a patch makes rather than its lines, would convert the patch-volume indicator into a measure that more closely tracks maintenance effort. It would also remove the dependence on the source-package format that confines the present metric to the quilt era.

A second direction broadens the comparison. Applying the same pipeline to other long-lived distributions, such as the Ubuntu and Arch systems studied by Horváth~\cite{horvath2012linuxdeps} or the Fedora system studied by Lin et al.~\cite{lin2022upstream}, would allow the trends reported here to be separated into those specific to Debian and those common to the ecosystem, and extending the analysis from the minbase variant to the standard and desktop tiers would test whether the accretion pattern observed in the minimal base also holds for larger installations. The pipeline could likewise incorporate the additional per-package metrics considered during scoping, in particular the Maintainability Index~\cite{Coleman1994Using}, alongside measures of coupling and cognitive complexity that the present control-flow metric does not capture.

A third direction takes up the inferential questions that the descriptive design of this thesis deliberately set aside. The co-visibility of the \texttt{systemd} transition across three dimensions, and the divergence between rising size and stationary per-function complexity, are observations that invite the correlational and causal analysis that larger samples and controlled comparisons would support. The per-release artefacts produced here are intended as the baseline for such work, and for the evaluation of automated software de-bloating techniques~\cite{quach2018debloating} against a documented historical record of how the base system has actually grown.
